%============================================================
% Part 2 opener  (generated by tools/outline.py; design: preamble/partpage.tex)
%============================================================
\dppartopener{Shortest Paths, Spanning Trees \& Modelling}{CF 1400 $\to$ CF 2200}{primaryblue}%
  {Weights enter the picture: shortest paths of every kind, minimum spanning trees, and the art of building the right graph.}

\dppartstats{6}{30}{60}{180}

\begin{dpjourney}
\dpstep{7}{Dijkstra's Algorithm}{Single-source shortest paths $\cdot$ Counting and tie-breaking $\cdot$ State-space Dijkstra $\cdot$ Several sources and reversed graphs $\cdot$ k shortest paths}{CF 1400 -- 2200}{no}
\dpstep{8}{Negative Weights, Bellman--Ford \& 0-1 BFS}{0-1 BFS $\cdot$ At most k edges $\cdot$ Negative cycles $\cdot$ Difference constraints $\cdot$ Mean and ratio cycles}{CF 1500 -- 2200}{no}
\dpstep{9}{All-Pairs Shortest Paths}{Floyd--Warshall $\cdot$ Vertex insertion order $\cdot$ Transitive closure $\cdot$ Minimax paths $\cdot$ Reconstructing from distances}{CF 1500 -- 2100}{no}
\dpstep{10}{Minimum Spanning Trees}{Kruskal and Prim $\cdot$ Virtual vertices $\cdot$ Which edges can be in an MST? $\cdot$ Bottleneck and minimax $\cdot$ Special and directed MSTs}{CF 1400 -- 2200}{no}
\dpstep{11}{Binary Lifting \& Jump Pointers}{k-th successor $\cdot$ Distances in functional graphs $\cdot$ Permutation powers and cycles $\cdot$ Doubling with aggregates $\cdot$ Jump pointers on arrays}{CF 1500 -- 2200}{no}
\dpstep{12}{Graph Construction Techniques}{Virtual vertices $\cdot$ Layered graphs $\cdot$ Segment-tree graphs $\cdot$ Complements and value vertices $\cdot$ Games on graphs (retrograde analysis)}{CF 1700 -- 2400}{no}
\end{dpjourney}

\dppartgoals{
  \item choose between BFS, 0-1 BFS, Dijkstra, Bellman--Ford and Floyd--Warshall from the constraints alone;
  \item prove and use the cut and cycle properties of minimum spanning trees;
  \item answer jump queries on functional graphs with binary lifting;
  \item build helper graphs: virtual vertices, layered states, segment-tree edges, complement traversal.
}
